% Faithful formal-publication transcription of complete Section4.
\section*{4 The construction}
\subsection*{4.1 Black-box error reduction}
Let $G:\{0,1\}^{s_0}\to\Sigma^n$ be an $(n,\Sigma,w,\epsilon_G)$ and $G_{\mathrm{aux}}:\{0,1\}^{s_{\mathrm{aux}}}\to(\{0,1\}^{s_0}\times\Sigma)^{k+1}$ be a $(k+1,\{0,1\}^{s_0}\times\Sigma,w,\epsilon_{\mathrm{aux}})$ PRG. Also, for $t\in\{0,1\}$ and $\sigma\in\Sigma$ we let
\begin{equation}\tag{9}G_{t,\ell}(x,\sigma)=\begin{cases}\sigma\circ G_{\ell-1}(x)&t=0,\\G_\ell(x)&t=1.\end{cases}\end{equation}
We now define the WPRG $(I,\mu):\{0,1\}^s\to\Sigma\times\mathbb R$. The seed $x\in\{0,1\}^s$ to our WPRG is interpreted as follows.
\begin{itemize}
\item The first $\log(k+1)$ bits encode $i\in\{0,\ldots,k\}$.
\item The next $\log\binom ni$ bits encode a sequence $\ell=(\ell_0,\ell_1,\ldots,\ell_i)$ such that $\ell_0+\cdots+\ell_i=n$, $\ell_i,\ldots,\ell_1>0$, and $\ell_0\ge0$.
\item The next $i$ bits are denoted by $t=t=(t_1,\ldots,t_i)\in\{0,1\}^i$.
\item The next $s_{\mathrm{aux}}$ bits are denoted by $x_{\mathrm{aux}}\in\{0,1\}^{s_{\mathrm{aux}}}$.
\end{itemize}
Overall, we can write $x=(i,\ell,t,x_{\mathrm{aux}})$, and the WPRG $(I,\mu)$ has seed length
\begin{equation}\tag{10}s=s_{\mathrm{aux}}+O(k\log n).\end{equation}
For brevity we sometimes omit the dependence of $i$, $(\ell_0,\ldots,\ell_i)$, $(t_1,\ldots,t_i)$, and $x_{\mathrm{aux}}$ on $x$. We define $I$ and $\mu$ as follows.
\[I(x)=\begin{cases}G_n(G_{\mathrm{aux}}(x_{\mathrm{aux}})_0)&i=0,\\G_{t_i,\ell_i}(G_{\mathrm{aux}}(x_{\mathrm{aux}})_i)\circ\cdots\circ G_{t_1,\ell_1}(G_{\mathrm{aux}}(x_{\mathrm{aux}})_1)\circ G_{\ell_0}(G_{\mathrm{aux}}(x_{\mathrm{aux}})_0)&\text{otherwise.}\end{cases}\]
\begin{equation}\tag{11}\mu(x)=\begin{cases}k+1&i=0,\\(k+1)\cdot\binom ni\cdot2^i\cdot(-1)^{t_1+\cdots+t_i}&\text{otherwise.}\end{cases}\end{equation}
where $G_{\mathrm{aux}}(x_{\mathrm{aux}})_j$ denotes the $j$'th symbol in $G_{\mathrm{aux}}(x_{\mathrm{aux}})\in(\{0,1\}^{s_0}\times\Sigma)^{k+1}$.
The weights are chosen so that the approximation yielded by the above WPRG is a derandomized version of Equation(8) for $(a,b)=(n+1,1)$. Note that in Equation(8) we used $r_1,\ldots,r_i$ which partitioned the interval $[n+1,1]$, while in Equation(11) we used $\ell_0,\ldots,\ell_i$ that sum to $n$. This is merely an alternative way of writing the sum -- the $\ell_i$-s are the sum of differences of the $r_i$-s.
\subsection*{4.2 Correctness}
In this section we use the same notation as in Section3.
\begin{authorclaim}{Lemma13}
Let $0<\epsilon<\epsilon_0=1/(4n)$ and let $k=\log_{1/\epsilon_0}(1/\epsilon)$. Suppose
\begin{itemize}
\item $G:\{0,1\}^{s_0}\to\Sigma^n$ is an $(n,\Sigma,w,\epsilon_G=\epsilon_0/(2(n+1)))$ PRG, and,
\item $G_{\mathrm{aux}}:\{0,1\}^{s_{\mathrm{aux}}}\to(\{0,1\}^{s_0}\times\Sigma)^{k+1}$ is a $(k+1,\{0,1\}^{s_0}\times\Sigma,w,\epsilon_{\mathrm{aux}}=\epsilon^3)$ PRG.
\end{itemize}
Then, $(I,\mu)$ is an $(n,\Sigma,w,\epsilon)$ WPRG with seed length $s=s_{\mathrm{aux}}+O(\log(1/\epsilon))$ computable in space $O(\operatorname{space}(G_{\mathrm{aux}})+\operatorname{space}(G)+\log s)$.
\end{authorclaim}
\begin{proof}
Assume $k,G$ and $G_{\mathrm{aux}}$ are as in the hypothesis of the lemma. The space complexity follows from Theorem7 and the seed length was analyzed in Equation(10). We are left to prove that $(I,\mu)$ is an $(n,\Sigma,w,\epsilon)$ WPRG. Fix any $(n,\Sigma,w)$ BP $B=(B_1,\ldots,B_n)$. Let $A$ be the $(n+1)w\times(n+1)w$ lower triangular block matrix in which
\[A[a,b]=\mathbb E_{x\in\{0,1\}^{s_0}}[\operatorname{val}(B_{[b,a-1]},G_{a-b}(x))]\]
for $a>b$, and $A[a,a]=I_w$. Since $G$ is $(n,\Sigma,w,\epsilon_G=\epsilon_0/(2(n+1)))$ PRG we have that $\|L^{-1}[a,b]-A[a,b]\|\le\epsilon_G$ and $\|L^{-1}-A\|\le(n+1)\epsilon_G$. By our choice of $\mu$,
\[\mathbb E_{x\in\{0,1\}^s}[\mu(x)\cdot\operatorname{val}(B,I(x))]=\sum_{i=0}^k\sum_{t,\ell}(-1)^{t_1+\cdots+t_i}\cdot\mathbb E_{x_{\mathrm{aux}}}[\operatorname{val}(B,I(i,\ell,t,x_{\mathrm{aux}}))],\]
and
\begin{align*}
R(A,L,k)[n+1,1]&=A[n+1,1]+\sum_{i=1}^k\sum_{t,r}(-1)^{t_1+\cdots+t_i}\cdot C_{t_i}[n+1,r_i]\cdots C_{t_1}[r_2,r_1]\cdot A[r_1,1],
\end{align*}
where $\ell_0+\cdots+\ell_i=n$ and $n+1>r_i>\cdots>r_1\ge1$. We soon prove:
\begin{authorclaim}{Claim14}
For every fixed $i\in\{0,\ldots,k\}$, $t\in\{0,1\}^i$, and $\ell$ such that $\ell_0+\cdots+\ell_i=n$
\[\left\|\mathbb E_{x_{\mathrm{aux}}}[\operatorname{val}(B,I(i,\ell,t,x_{\mathrm{aux}}))]-C_{t_i}[n+1,r_i]\cdots C_{t_1}[r_2,r_1]\cdot A[r_1,1]\right\|\le\epsilon_{\mathrm{aux}},\]
where $r_j=1+\ell_0+\cdots+\ell_{j-1}$.
\end{authorclaim}
As we have at most $(k+1)n^k2^k$ summands, we see that
\begin{align*}
\left\|\mathbb E_{x\in\{0,1\}^s}[\mu(x)\cdot\operatorname{val}(B,I(x))]-R(A,L,k)[n+1,1]\right\|
&\le(k+1)n^k2^k\cdot\epsilon_{\mathrm{aux}}\\
&\le\frac{n^{2k}}2\cdot\epsilon_{\mathrm{aux}}\le\frac\epsilon2.
\end{align*}
It therefore follows from Claim12 that
\begin{align*}
\left\|R(A,L,k)[n+1,1]-\mathbb E_{x\in\Sigma^n}[\operatorname{val}(B,x)]\right\|
&\le(n+1)(2(n+1)\epsilon_G)^{k+1}\\
&\le2n\cdot\epsilon_0^{k+1}\le2n\epsilon_0\epsilon=\frac\epsilon2,
\end{align*}
which together completes the proof.
\end{proof}
\begin{proof}[Proof of Claim14]
Fix $i\in\{0,\ldots,k\}$, $\ell_0+\cdots+\ell_i=n$, and $t\in\{0,1\}^i$ and recall that $r_j=1+\ell_0+\cdots+\ell_{j-1}$. We define a $(k+1,\{0,1\}^{s_0}\times\Sigma,w)$ BP $B'=(B'_0,\ldots,B'_k)$ (that depends on $i,\ell$, and $t$) such that for all $j=0,\ldots,k$,
\begin{equation}\tag{12}B'_j(x,\sigma)=\begin{cases}\operatorname{val}(B_{[r_j,r_{j+1}-1]},\sigma\circ G_{\ell_j-1}(x))&j>0,t=0,\\\operatorname{val}(B_{[r_j,r_{j+1}-1]},G_{\ell_j}(x))&j>0,t=1,\\\operatorname{val}(B_{[1,r_1-1]},G_{\ell_0}(x))&j=0.\end{cases}\end{equation}
We stress that $B'_j$ is a BP because a product of Boolean stochastic matrices is Boolean stochastic. The claim now follows since $G_{\mathrm{aux}}$ is a $(k+1,\{0,1\}^{s_0}\times\Sigma,w,\epsilon_{\mathrm{aux}})$ PRG.
\end{proof}
\subsection*{4.3 The final construction}
We now instantiate Lemma13 with $G_{\mathrm{aux}}$ being the INW PRG from Theorem9 and $G$ being an arbitrary PRG. The reason for using the INW generator is its additive dependence on $\log|\Sigma|$.
\begin{authorclaim}{Corollary15}
Let $G:\{0,1\}^{s_0}\to\Sigma^n$ be an $(n,\Sigma,w,\epsilon_G)$. Then, for any error parameter $1/(4n)>\epsilon>0$ there exists an $(n,\Sigma,w,\epsilon)$ WPRG with seed length
\[s_0+O\left(\log\frac w\epsilon\cdot\log\log_n\frac1\epsilon\right)\]
computable in space $O(\operatorname{space}(G)+\log\log_n(1/\epsilon)\cdot(\log\log(w/\epsilon))^2)$.
\end{authorclaim}
Had we used Nisan's PRG from Theorem8 instead of INW then the seed length would deteriorate to
\[O\left(s_0\cdot\log\log_n\frac1\epsilon+\log\frac w\epsilon\cdot\log\log_n\frac1\epsilon\right).\]
Corollary15 can be interpreted as an error reduction procedure for PRGs with a slight overhead in the seed and space complexity. We proceed by applying this error reduction to Nisan's PRG from Theorem8.
\begin{authorclaim}{Corollary16}
For any positive integers $n,w$, any error parameter $1/(4n)>\epsilon>0$ and any alphabet $\Sigma$, there exists an $(n,\Sigma,w,\epsilon)$ WPRG with seed length
\[O\left(\log n\log(nw|\Sigma|)+\log\frac w\epsilon\cdot\log\log_n\frac1\epsilon\right)\]
computable in space $O(\log(nw|\Sigma|)+\log\log_n(1/\epsilon)\cdot(\log\log(w/\epsilon))^2)$.
\end{authorclaim}
Note that for $\epsilon$ which is not tiny the space complexity is dominated by the first term. Specifically, for $\epsilon>2^{-2^{\log^{1/3}n}}$, $w<2^{2^{\log^{1/3}n}}$ the space complexity is indeed $O(\log(nw|\Sigma|))$. Had we used INW instead, the space complexity would deteriorate to
\[O\left(\log n\cdot\left[\log\log\frac{nw|\Sigma|}\epsilon\right]^2+\log\frac w\epsilon\cdot\log\log_n\frac1\epsilon\right).\]
